% ============================================================
\section{Split Transform, Schur Complement, and Nonlinear Elimination}
\label{sec:resolvent-defect}
% ============================================================

\begin{quote}\small
\textbf{What this automates.}
When inverting or resolving an operator under an orthogonal projection $P$, the ``shadow'' degrees of freedom (in $\CCComp\cH$, where $\CCComp = I - P$) contribute a correction term to the observed resolvent.
This section provides the exact Feshbach--Schur identity that computes this
correction, namely the \emph{Schur correction} (or \emph{self-energy}),
without approximation or series expansion.
It then upgrades the same split algebra to nonlinear equations by exact shadow
elimination on a graph.

\smallskip\noindent
\textbf{Generality.}
The constructions apply to \emph{any} orthogonal projection $P$, not just horizon projections $\CCPi{h}$.
The horizon case is recovered by setting $P = \CCPi{h}$.
Field-induced horizons (Section~\ref{subsec:field-induced-horizons}) give $P = \CCSpProj{\CCField}{S}$.
\end{quote}

\subsection{Block decomposition of operators}

\begin{definition}[P/Q splitting]
\label{def:pq-split}\lean{PQSplitting}
Let $P$ be an orthogonal projection on $\cH$ with complement $\CCComp := I - P$.
Define the subspaces:
\[
\cH_P := P \cH, \qquad \cH_{\CCComp} := \CCComp \cH.
\]
Then $\cH = \cH_P \oplus \cH_{\CCComp}$ as an orthogonal direct sum.
\end{definition}

\begin{definition}[Block components of an operator]
\label{def:block-components}\lean{blockPP}\lean{blockPQ}\lean{blockQP}\lean{blockQQ}
For a linear operator $A:\cH \to \cH$, define the block components:
\begin{align*}
\CCBlockPP{A} &:= P A P &&\text{($P \to P$ block)}, \\
\CCBlockPQ{A} &:= P A \CCComp &&\text{($\CCComp \to P$ block)}, \\
\CCBlockQP{A} &:= \CCComp A P &&\text{($P \to \CCComp$ block)}, \\
\CCBlockQQ{A} &:= \CCComp A \CCComp &&\text{($\CCComp \to \CCComp$ block)}.
\end{align*}
In block form relative to $\cH = \cH_P \oplus \cH_{\CCComp}$:
\[
A = \begin{pmatrix} \CCBlockPP{A} & \CCBlockPQ{A} \\ \CCBlockQP{A} & \CCBlockQQ{A} \end{pmatrix}.
\]
\end{definition}

\paragraph{Connection to horizon notation.}
When $P = \CCPi{h}$ is a horizon projection, the block notation corresponds to:
\begin{align*}
\CCBlockPP{A} &= A_{oo} = \CCRestr{A}{h}, &
\CCBlockPQ{A} &= A_{os} = \CCPi{h} A \CCPibar{h}, \\
\CCBlockQP{A} &= A_{so} = \CCPibar{h} A \CCPi{h} = \CCLeak{A}{h}, &
\CCBlockQQ{A} &= A_{ss} = \CCPibar{h} A \CCPibar{h}.
\end{align*}

\subsection{The split transform (forward direction)}

\begin{definition}[Resolvent]
\label{def:resolvent}\lean{HasResolvent}
For a linear operator $A$ on $\cH$ and $z \in \mathbb{C}$ such that $(I - zA)$ is invertible, define the \emph{resolvent}
\[
R_A(z) := (I - zA)^{-1}.
\]
\end{definition}

\begin{definition}[Schur correction / self-energy]
\label{def:schur-correction}\lean{schurComplement}
For an orthogonal projection $P$ with $\CCComp = I - P$, define the \emph{Schur correction} (or \emph{self-energy}) as
\[
\CCSigmaP{P}{z} := \CCBlockPQ{A} (I - z \CCBlockQQ{A})^{-1} \CCBlockQP{A},
\]
defined whenever $(I - z \CCBlockQQ{A})$ is invertible on $\cH_{\CCComp}$.
\end{definition}

\begin{definition}[Effective operator]
\label{def:effective-operator}\lean{effectiveOp}
The \emph{effective operator} (or \emph{observed effective operator}) is
\[
\CCAeff{A}{P}(z) := \CCBlockPP{A} + z \CCSigmaP{P}{z}.
\]
\end{definition}

\begin{definition}[Split transform]
\label{def:split-transform}\lean{splitTransform}
The \emph{split transform} is the map
\[
\CCSplit{P}{A}{z} := \bigl( \CCAeff{A}{P}(z),\, \CCSigmaP{P}{z} \bigr)
\]
that takes $(A, P, z)$ and returns the effective operator and Schur correction.
\end{definition}

\begin{theorem}[Forward split / Feshbach--Schur identity]
\label{thm:feshbach-schur}\lean{schur_factorization}
Let $A:\cH \to \cH$ be a bounded linear operator, $P$ an orthogonal projection, and $z \in \mathbb{C}$ such that both $(I - zA)$ on $\cH$ and $(I - z \CCBlockQQ{A})$ on $\cH_{\CCComp}$ are invertible.
Then the observed resolvent satisfies
\[
P (I - zA)^{-1} P
\;=\;
\bigl( I - z \CCAeff{A}{P}(z) \bigr)^{-1}
\;=\;
\bigl( I - z \CCBlockPP{A} - z^2 \CCSigmaP{P}{z} \bigr)^{-1}
\]
as operators on $\cH_P$.
\end{theorem}

(Proof in Appendix~\ref{sec:appendix-elimination}.)

This is the classical Schur-complement / Feshbach elimination pattern
\cite{Haynsworth1968,Zhang2005,Feshbach1958,Feshbach1962}, fixed here in the
notation used throughout the horizon calculus.

% --- Schur Complement Diagram ---
\begin{figure}[ht]
\centering
\resizebox{0.78\textwidth}{!}{%
\begin{tikzpicture}[node distance=3.0cm]
  \node[ccbox, minimum width=2.4cm] (p) at (-2.8,0.8) {$\cH_P$};
  \node[ccbox, minimum width=2.6cm] (q) at (2.8,0.8) {$\cH_{\CCComp}$};
  \node[ccbox, align=center, minimum width=5.8cm] (sigma) at (0,-2.0)
    {$\CCSigmaP{P}{z}=\CCBlockPQ{A}(I-z\CCBlockQQ{A})^{-1}\CCBlockQP{A}$};

  \draw[ccarrow, bend left=18] (p.north east)
    to node[above, font=\scriptsize, pos=0.52] {$\CCBlockQP{A}$} (q.north west);
  \draw[ccarrow, bend left=18] (q.south west)
    to node[below, font=\scriptsize, pos=0.52] {$\CCBlockPQ{A}$} (p.south east);
  \draw[ccdarrow] (q) to[loop right, looseness=4]
    node[right, font=\scriptsize] {$(I-z\CCBlockQQ{A})^{-1}$} (q);
  \draw[ccarrow] ($(p)!0.5!(q)$) -- node[right, font=\scriptsize]
    {round trip through $\CCComp$} (sigma);
\end{tikzpicture}%
}
\caption{Schur complement / split transform: round-trip through $\CCComp$-block produces the Schur correction $\CCSigmaP{P}{z}$.}
\label{fig:schur-complement}
\end{figure}

\begin{quote}\small
\textbf{Intuition.}
The Schur correction $\CCSigmaP{P}{z}$ encodes how the $\CCComp$-degrees of freedom modify the effective resolvent.
$P$-states leak into $\CCComp$ ($\CCBlockQP{A}$), propagate there ($(I - z \CCBlockQQ{A})^{-1}$), and return ($\CCBlockPQ{A}$).
This round-trip contribution corrects the observed resolvent beyond the naive projection.
\end{quote}

\paragraph{Exact resummation.}
The identity is exact; no truncation or perturbative expansion is required.
The Schur correction $\CCSigmaP{P}{z}$ resums all paths through $\cH_{\CCComp}$ into a single correction term.
This is the Coherence Calculus mechanism for ``integrating out'' degrees of freedom.
The exact coherence and recovery theorems below show that, on a faithful tower,
this eliminated contribution remains part of the same underlying operator rather
than an added constitutive source.

\subsection{Residual signature (transport characterization)}

\begin{definition}[Transport blocks and residual signature]
\label{def:residual-signature}\lean{ResidualSignature}
For an operator $A$ and projection $P$, define:
\begin{itemize}[itemsep=0.2em]
\item \textbf{Transport out:} $\CCTransOut{P} := \CCBlockQP{A}$ (the $P \to \CCComp$ block).
\item \textbf{Transport in:} $\CCTransIn{P} := \CCBlockPQ{A}$ (the $\CCComp \to P$ block).
\item \textbf{Internal block:} $\CCBlockQQ{A}$ (the $\CCComp \to \CCComp$ block).
\end{itemize}
The \emph{residual signature} of $A$ with respect to $P$ is the triple $(\CCTransIn{P}, \CCBlockQQ{A}, \CCTransOut{P})$.
\end{definition}

\paragraph{Residual determines correction.}
The Schur correction is completely determined by the residual signature:
\[
\CCSigmaP{P}{z} = \CCTransIn{P} (I - z \CCBlockQQ{A})^{-1} \CCTransOut{P}.
\]
The round-trip factorization: transport out, propagate internally, transport back in.

\begin{definition}[Residual return field on a cut]
\label{def:residual-return-field}\lean{ResidualReturnField}
Fix an operator \(A\), a projection \(P\), and a shadow propagator
\[
R_Q=(I-z\CCBlockQQ{A})^{-1}
\]
whenever this inverse exists. The \emph{residual return field} generated by a
visible state \(u\in P\cH\) is the shadow datum
\[
\chi_{P,A,z}(u):=R_Q\,\CCTransOut{P}u.
\]
Its \emph{returned residual forcing} on the visible sector is
\[
\mathcal F^{\mathrm{return}}_{P,A,z}(u)
:=
\CCTransIn{P}\chi_{P,A,z}(u).
\]
\end{definition}

\begin{proposition}[Returned residual is exactly the Schur correction]
\label{prop:return-residual-is-schur}\lean{returnedResidual_eq_schur}
In the setting of Definition~\ref{def:residual-return-field},
\[
\mathcal F^{\mathrm{return}}_{P,A,z}(u)=\CCSigmaP{P}{z}u
\qquad (u\in P\cH).
\]
Thus the Schur correction is not an extra constitutive term. It is exactly the
visible trace of a residual field propagated in the eliminated sector and then
returned to observation.
\end{proposition}

(Proof in Appendix~\ref{sec:appendix-elimination}.)

\begin{corollary}[Self-adjoint structure]
\label{cor:self-adjoint-sigma}\lean{selfadjoint_block_relation}\lean{schur_selfadjoint_leading}
If $A$ is self-adjoint and $z \in \bbR$ with $(I - z\CCBlockQQ{A})^{-1}$ self-adjoint (holds when $z$ is not in the spectrum of $\CCBlockQQ{A}$ scaled appropriately), then:
\begin{enumerate}[label=(\roman*)]
\item $\CCTransIn{P} = (\CCTransOut{P})^*$, i.e., $\CCBlockPQ{A} = (\CCBlockQP{A})^*$.
\item $\CCSigmaP{P}{z}$ is self-adjoint on $\cH_P$.
\end{enumerate}
\end{corollary}

\begin{proof}
(i) For self-adjoint $A$: $(PAQ)^* = Q^* A^* P^* = QAP$, so $\CCBlockPQ{A}^* = \CCBlockQP{A}$.
(ii) $\CCSigmaP{P}{z}^* = (\CCBlockPQ{A} R_Q \CCBlockQP{A})^* = \CCBlockQP{A}^* R_Q^* \CCBlockPQ{A}^* = \CCBlockPQ{A} R_Q \CCBlockQP{A} = \CCSigmaP{P}{z}$.
\end{proof}

\subsection{Geometric series expansion}

\paragraph{Geometric series expansion.}
\label{rem:geometric-expansion}
If $\|z \CCBlockQQ{A}\| < 1$, then $(I - z \CCBlockQQ{A})^{-1} = \sum_{k=0}^\infty z^k (\CCBlockQQ{A})^k$ converges, giving
\[
\CCSigmaP{P}{z} = \sum_{k=0}^\infty z^k \CCBlockPQ{A} (\CCBlockQQ{A})^k \CCBlockQP{A}.
\]
This expands the Schur correction as a sum over ``$\CCComp$-excursions'' of increasing length.
The exact identity (Theorem~\ref{thm:feshbach-schur}) holds without this convergence assumption.

\paragraph{Connection to defect energy.}
When $P = \CCPi{h}$ and $A = d$ is a nilpotent differential, the leading term $\CCBlockPQ{A} \CCBlockQP{A} = \CCPi{h} d \CCPibar{h} \cdot \CCPibar{h} d \CCPi{h}$ is the defect operator $\CCDefect{d}{h}$ from the Horizon Fundamental Theorem (Section~\ref{sec:hft}).
The Schur correction generalizes this to all orders in $z$.

\subsection{Schur transport jets}
\label{subsec:schur-transport-jet}

\begin{definition}[Schur transport jet]
\label{def:schur-transport-jet}\lean{schur_transport_jet}
Let $(\CCPi{h})_{h\in\bbN}$ be a faithful horizon tower, let
$A:\cH\to\cH$ be bounded, and fix a horizon $h$. Assume
\[
\|z\,\CCPibar{h}A\CCPibar{h}\|<1.
\]
The \emph{Schur transport jet} at horizon $h$ is the family of operators
\[
K_{h,0}(A):=\CCPi{h}A\CCPi{h},
\qquad
K_{h,n}(A):=\CCPi{h}A(\CCPibar{h}A\CCPibar{h})^{n-1}\CCPibar{h}A\CCPi{h}
\quad(n\ge1),
\]
so that the effective operator admits the exact expansion
\[
\CCAeff{A}{\CCPi{h}}(z)=K_{h,0}(A)+\sum_{n\ge1} z^n K_{h,n}(A).
\]
\end{definition}

\begin{theorem}[Constructive order profile of the Schur transport jet]
\label{thm:schur-transport-jet-order}\lean{schur_transport_jet_order}
Assume the setting of Definition~\ref{def:schur-transport-jet}, and suppose the
underlying operator $A$ has tower transport order at most $r$ in the sense of
Definition~\ref{def:tower-transport-order}. Then:
\begin{enumerate}[label=(\roman*), leftmargin=2em, itemsep=0.2em]
\item $K_{h,0}(A)$ has order at most $r$;
\item for every $n\ge1$, the coefficient $K_{h,n}(A)$ has order at most
      $(n+1)r$;
\item the exact effective operator $\CCAeff{A}{\CCPi{h}}(z)$ therefore carries a
      constructive finite-order jet at every horizon, with order bounds
      determined solely by the tower order of $A$.
\end{enumerate}
\end{theorem}

(Proof in Appendix~\ref{sec:appendix-elimination}.)

\subsection{Stability bounds}

\begin{proposition}[Schur correction bound]
\label{prop:sigma-bound}\lean{schur_factorization_bound}
Assume $(I - z\CCBlockQQ{A})^{-1}$ exists with $\|(I - z\CCBlockQQ{A})^{-1}\| \le M$.
Then
\[
\|\CCSigmaP{P}{z}\| \;\le\; \|\CCBlockPQ{A}\| \cdot M \cdot \|\CCBlockQP{A}\|.
\]
\end{proposition}

\begin{proof}
$\|\CCSigmaP{P}{z}\| = \|\CCBlockPQ{A} (I - z\CCBlockQQ{A})^{-1} \CCBlockQP{A}\| \le \|\CCBlockPQ{A}\| \cdot \|(I - z\CCBlockQQ{A})^{-1}\| \cdot \|\CCBlockQP{A}\|$.
\end{proof}

\begin{proposition}[Perturbation stability]
\label{prop:perturbation-stability}\lean{PerturbationData}\lean{schurPerturbationDiff}
Let $A, \tilde{A}$ be operators with $E := \tilde{A} - A$.
Assume:
\begin{itemize}[itemsep=0.2em]
\item $(I - z\CCBlockQQ{A})^{-1}$ exists with $\|(I - z\CCBlockQQ{A})^{-1}\| \le M$.
\item $\|z \CCBlockQQ{E}\| \cdot M < 1$.
\end{itemize}
Then $(I - z\CCBlockQQ{\tilde{A}})^{-1}$ exists and
\[
\|\CCSigmaP{P}{z}[\tilde{A}] - \CCSigmaP{P}{z}[A]\|
\;\le\;
C(M, \|A\|, |z|) \cdot \|E\|,
\]
where
\[
C(M, \|A\|, |z|) = M^2 \bigl( 2 + |z| M (\|\CCBlockQQ{A}\| + \|E\|) \bigr) + M.
\]
\end{proposition}

\begin{proof}
Use the resolvent identity $(I - z\CCBlockQQ{\tilde{A}})^{-1} - (I - z\CCBlockQQ{A})^{-1} = z(I - z\CCBlockQQ{\tilde{A}})^{-1} \CCBlockQQ{E} (I - z\CCBlockQQ{A})^{-1}$.
The Neumann series gives $\|(I - z\CCBlockQQ{\tilde{A}})^{-1}\| \le M/(1 - |z| M \|\CCBlockQQ{E}\|)$.
Expand the difference $\CCSigmaP{P}{z}[\tilde{A}] - \CCSigmaP{P}{z}[A]$ using the product rule and collect terms.
\end{proof}

\subsection{Nested splits (band-by-band elimination)}

\begin{lemma}[Associativity of Schur complement]
\label{lem:nested-schur}\lean{NestedProjections}\lean{nestedSchurAssociativity}
Let $P_1, P_2$ be orthogonal projections with $P_1 \le P_2$ (i.e., $P_1 P_2 = P_1$).
Define $\CCComp_1 = I - P_1$ and $\CCComp_2 = I - P_2$, so $\CCComp_2 \le \CCComp_1$.
Under appropriate invertibility conditions, eliminating $\CCComp_1$ in two stages (first eliminate $\CCComp_2 \setminus P_2$, then eliminate $(P_2 \setminus P_1)$) equals direct elimination of $\CCComp_1$:
\[
\CCSigmaP{P_1}{z}[A] = \CCSigmaP{P_1}{z}[\CCAeff{A}{P_2}(z)].
\]
\end{lemma}

\begin{proof}
This follows from the associativity of Schur complements in block matrix theory: if $M$ is a $3 \times 3$ block matrix, the Schur complement onto the first block can be computed by first eliminating the third block, then the second.
The invertibility conditions transfer through the chain.
\end{proof}

\paragraph{Horizon-to-horizon elimination.}
When $P_1 = \CCPi{h}$ and $P_2 = \CCPi{h'}$ with $h < h'$, Lemma~\ref{lem:nested-schur} says we can eliminate from $h'$ to $h$ in one step, or go via intermediate levels.
This justifies recursive/hierarchical elimination schemes.

\subsection{Exact Schur towers and faithful recovery}
\label{subsec:exact-schur-towers}

\begin{definition}[Exact Schur tower at fixed parameter]
\label{def:exact-schur-tower-general}\lean{exact_schur_tower_general}
Let $(\CCPi{h})_{h\in\bbN}$ be a faithful horizon tower on a Hilbert space
$\cH$, let $A:\cH\to\cH$ be bounded, and fix
$z\in\mathbb{C}$ such that, for every $h$, the shadow resolvent
\[
(I-z\CCPibar{h}A\CCPibar{h})^{-1}
\]
exists. The associated \emph{exact Schur tower} is the family of effective
operators
\[
A_{\mathrm{eff},h}(z)
:=
\CCPi{h}A\CCPi{h}
+
z\,\CCPi{h}A\CCPibar{h}(I-z\CCPibar{h}A\CCPibar{h})^{-1}\CCPibar{h}A\CCPi{h}
=
\CCAeff{A}{\CCPi{h}}(z).
\]
\end{definition}

\begin{definition}[Schur coherence across nested horizons]
\label{def:schur-coherent-family-general}\lean{schur_coherent_family_general}
Let $\{A_{\mathrm{eff},h}(z)\}_{h\in\bbN}$ be the exact Schur tower of
Definition~\ref{def:exact-schur-tower-general}. For $h\le k$, let
\[
P_{h\mid k}:=\CCPi{h}\big|_{\CCPi{k}\cH}
\]
denote the induced orthogonal projection from the $k$th observed sector onto
the $h$th one. The family is \emph{Schur-coherent} if
\[
A_{\mathrm{eff},h}(z)
=
\CCAeff{A_{\mathrm{eff},k}(z)}{P_{h\mid k}}(z)
\qquad (h\le k).
\]
Thus cross-horizon compatibility is expressed by exact nested elimination,
rather than by literal intertwining of operators between different observed
sectors.
\end{definition}

\begin{theorem}[Exact Schur coherence of a bounded-operator tower]
\label{thm:exact-schur-coherence-general}\lean{exact_schur_coherence_general}
Let $(\CCPi{h})_{h\in\bbN}$ be a faithful horizon tower on $\cH$, let
$A:\cH\to\cH$ be bounded, and fix
$z\in\mathbb{C}$ such that every shadow resolvent
\[
(I-z\CCPibar{h}A\CCPibar{h})^{-1}
\]
exists. Then the exact Schur tower of
Definition~\ref{def:exact-schur-tower-general} is Schur-coherent in the sense
of Definition~\ref{def:schur-coherent-family-general}. Equivalently, for every
$h\le k$,
\[
\CCAeff{A}{\CCPi{h}}(z)
=
\CCAeff{\CCAeff{A}{\CCPi{k}}(z)}{P_{h\mid k}}(z).
\]
\end{theorem}

(Proof in Appendix~\ref{sec:appendix-elimination}.)

\begin{theorem}[Strong recovery of the underlying operator from the exact Schur tower]
\label{thm:schur-tower-strong-recovery-general}\lean{schur_tower_strong_recovery_general}
Let $(\CCPi{h})_{h\in\bbN}$ be a faithful horizon tower on a Hilbert space
$\cH$, let $A:\cH\to\cH$ be bounded, and fix
$z\in\mathbb{C}$ such that every shadow resolvent exists and satisfies the
uniform bound
\[
\sup_{h\in\bbN}
\left\|
(I-z\CCPibar{h}A\CCPibar{h})^{-1}
\right\|
\le M
\]
for some $M\ge 0$. Then for every $x\in\cH$,
\[
A_{\mathrm{eff},h}(z)x \to Ax
\qquad\text{as }h\to\infty,
\]
where $A_{\mathrm{eff},h}(z)$ is the exact Schur tower of
Definition~\ref{def:exact-schur-tower-general}.
\end{theorem}

(Proof in Appendix~\ref{sec:appendix-elimination}.)

\begin{remark}[Faithful recovery of the residual return field]
\label{rem:faithful-recovery-residual-return}
\lean{exact_schur_coherence_general, schur_tower_strong_recovery_general, nonlinear_shadow_elimination}
Together, Theorems~\ref{thm:exact-schur-coherence-general}
and \ref{thm:schur-tower-strong-recovery-general} show that a faithful exact
Schur tower recovers the underlying operator itself from the effective laws
seen at finite horizons. The residual return field of
Definition~\ref{def:residual-return-field} is therefore not a free source
insertion. Combined with exact nonlinear shadow elimination
(Theorem~\ref{thm:nonlinear-shadow-elimination}), this means that the residual
return field is an eliminated hidden datum whose returned visible trace is
exactly controlled by the same observed effective law; it is not an independent
postulate layered on top of that law.
\end{remark}

\begin{theorem}[Vanishing-normalized selected Schur families recover the same continuum operator]
\label{thm:selected-schur-recovery-general}\lean{selected_schur_recovery_general}
Let $(\CCPi{h})_{h\in\bbN}$ be a faithful horizon tower on $\cH$, let
$A:\cH\to\cH$ be bounded, and fix
$z\in\mathbb{C}$ such that the hypotheses of
Theorem~\ref{thm:schur-tower-strong-recovery-general} hold. For each horizon
$h$, let
\[
C_h:\CCPi{h}\cH\to\CCPi{h}\cH
\]
be any selected correction. Define the selected effective family by
\[
\widetilde A_h(z):=A_{\mathrm{eff},h}(z)+C_h.
\]
If
\[
\|C_h\CCPi{h}x\|\to 0
\qquad\text{for every }x\in\cH,
\]
then
\[
\widetilde A_h(z)x\to Ax
\qquad\text{for every }x\in\cH.
\]
\end{theorem}

(Proof in Appendix~\ref{sec:appendix-elimination}.)

\subsection{Corollary for horizon case}

\begin{corollary}[Horizon specialization]
\label{cor:horizon-specialization}\lean{HorizonSpecialization}\lean{horizon_schur_via_spec}
Setting $P = \CCPi{h}$ and $\CCComp = \CCPibar{h}$, all results in this section apply with:
\begin{align*}
\CCBlockPP{A} &= \CCRestr{A}{h}, & \CCBlockPQ{A} &= \CCPi{h} A \CCPibar{h}, \\
\CCBlockQP{A} &= \CCLeak{A}{h}, & \CCBlockQQ{A} &= \CCPibar{h} A \CCPibar{h}.
\end{align*}
The Schur correction becomes $\CCSigma{h}{z}$ as in the original notation.
\end{corollary}

\paragraph{Scale variation of resolvent.}
The horizon flow $\CCFlow{\cdot}{h}$ (Section~\ref{sec:advanced-operators}) applied to the observed resolvent $\CCPi{h} (I-zA)^{-1} \CCPi{h}$ provides a canonical scale variation operator for studying how the effective resolvent changes across coherence levels.

\subsection{Determinant factorization (finite-dimensional case)}

\begin{theorem}[Determinant factorization]
\label{thm:det-factorization}\lean{CoherenceCalculus.det_factorization}
Let $A:\cH \to \cH$ be a linear operator on a finite-dimensional Hilbert space with orthogonal projection $P$ and $\CCComp = I - P$.
For $z \in \mathbb{C}$ such that $(I - z\CCBlockQQ{A})$ is invertible, the characteristic determinant factors as:
\[
\det(I - zA) = \det(I - z\CCBlockQQ{A}) \cdot \det\bigl(I - z\CCAeff{A}{P}(z)\bigr),
\]
where $\CCAeff{A}{P}(z) = \CCBlockPP{A} + z\CCSigmaP{P}{z}$ is the effective operator (Definition~\ref{def:effective-operator}).
\end{theorem}

(Proof in Appendix~\ref{sec:appendix-elimination}.)

\begin{corollary}[Determinant ratio identity]
\label{cor:det-ratio}\lean{CoherenceCalculus.det_ratio}
Under the hypotheses of Theorem~\ref{thm:det-factorization}:
\[
\frac{\det(I - zA)}{\det(I - z\CCBlockQQ{A})} = \det\bigl(I - z\CCAeff{A}{P}(z)\bigr).
\]
This expresses the ``observable determinant ratio'' purely in terms of the effective operator.
\end{corollary}

\paragraph{Spectral interpretation.}
The identity separates the characteristic polynomial of $A$ into contributions from the $\CCComp$-block and the $P$-effective operator.
The roots of the left-hand side (eigenvalues of $A$ scaled by $1/z$) decompose into those ``coming from'' the shadow block and those ``coming from'' the effective observed dynamics.

\subsection{Functoriality under unitary conjugation}

\begin{theorem}[Conjugacy invariance of split transform]
\label{thm:split-functoriality}\lean{CoherenceCalculus.split_functoriality}
Let $U:\cH \to \cH$ be unitary, and define $P' := UPU^*$ and $A' := UAU^*$.
Then:
\begin{enumerate}[label=(\roman*)]
\item $\CCSigmaP{P'}{z}[A'] = U \CCSigmaP{P}{z}[A] U^*$.
\item $\CCAeff{A'}{P'}(z) = U \CCAeff{A}{P}(z) U^*$.
\end{enumerate}
In other words, the split transform commutes with unitary conjugation.
\end{theorem}

(Proof in Appendix~\ref{sec:appendix-elimination}.)

\paragraph{Symmetry inheritance.}
Theorem~\ref{thm:split-functoriality} implies that if a symmetry group $G$ acts unitarily on $\cH$ and preserves both $A$ and $P$ (i.e., $UAU^* = A$ and $UPU^* = P$ for all $U \in G$), then $G$ also preserves the Schur correction $\CCSigmaP{P}{z}$ and effective operator $\CCAeff{A}{P}(z)$.
This is the mechanism by which symmetry constraints propagate through the split transform.

\subsection{Exact nonlinear shadow elimination}

The linear Schur identity eliminates the shadow sector at the level of
resolvents. For nonlinear equations the correct analog is different in form but
identical in purpose: solve the shadow equation exactly as a graph over the
observed variables, then substitute that graph back into the observed equation.
The result is an exact reduced law on the observed sector, not an ansatz.

\begin{definition}[Nonlinear $P/Q$ split]
\label{def:nonlinear-pq-split}\lean{nonlinear_pq_split}
Let $P$ be an orthogonal projection on $\cH$, set $Q:=I-P$, and let
$\mathcal D\subseteq \cH$ be the domain of a map
$F:\mathcal D\to \cH$. For $(p,q)\in \cH_P\times\cH_Q$ with
$p+q\in\mathcal D$, define
\[
F_P(p,q):=PF(p+q),
\qquad
F_Q(p,q):=QF(p+q).
\]
Thus the equation $F(u)=0$ with $u=p+q$ is equivalent to the coupled system
\[
F_P(p,q)=0,
\qquad
F_Q(p,q)=0.
\]
\end{definition}

\begin{definition}[Exact nonlinear effective law]
\label{def:nonlinear-effective-law}\lean{nonlinear_effective_law}
Let $U_P\subseteq \cH_P$ and $U_Q\subseteq \cH_Q$. Suppose
$\Psi:U_P\to U_Q$ is a map such that $p+\Psi(p)\in\mathcal D$ for every
$p\in U_P$. The associated \emph{effective nonlinear law} on $U_P$ is
\[
F_{\mathrm{eff}}(p):=F_P\bigl(p,\Psi(p)\bigr)
=
P F\bigl(p+\Psi(p)\bigr).
\]
\end{definition}

\begin{theorem}[Exact nonlinear shadow elimination]
\label{thm:nonlinear-shadow-elimination}\lean{nonlinear_shadow_elimination}
Let $F:\mathcal D\to\cH$ be a map and let $P$ be an orthogonal projection with
$Q=I-P$. Let $U_P\subseteq \cH_P$ and $U_Q\subseteq \cH_Q$, and suppose that a
map $\Psi:U_P\to U_Q$ satisfies:
\begin{enumerate}[label=(\roman*), leftmargin=2em, itemsep=0.2em]
\item for every $p\in U_P$, one has $p+\Psi(p)\in\mathcal D$ and
      $F_Q\bigl(p,\Psi(p)\bigr)=0$;
\item whenever $p\in U_P$, $q\in U_Q$, $p+q\in\mathcal D$, and
      $F_Q(p,q)=0$, then $q=\Psi(p)$.
\end{enumerate}
Then the map
\[
p \longmapsto p+\Psi(p)
\]
is a bijection between the reduced solution set
\[
\{\,p\in U_P : F_{\mathrm{eff}}(p)=0\,\}
\]
and the full solution set
\[
\{\,u=p+q\in\mathcal D : p\in U_P,\ q\in U_Q,\ F(u)=0\,\}.
\]
Equivalently, on $U_P\times U_Q$ the full nonlinear equation $F(u)=0$ is
exactly equivalent to the reduced observed equation
\[
F_{\mathrm{eff}}(p)=0
\]
together with the reconstruction rule $q=\Psi(p)$.
\end{theorem}

(Proof in Appendix~\ref{sec:appendix-elimination}.)

\begin{figure}[ht]
\centering
\resizebox{0.8\textwidth}{!}{%
\begin{tikzpicture}[node distance=2.6cm]
  \node[ccbox, minimum width=1.9cm] (p) at (-4.3,1.2) {$p \in \cH_P$};
  \node[ccbox, minimum width=2.2cm] (shadow) at (-0.9,1.2) {$F_Q(p,q)=0$};
  \node[ccbox, minimum width=2.0cm] (psi) at (2.5,1.2) {$q=\Psi(p)$};
  \node[ccbox, minimum width=2.5cm] (u) at (-0.9,-1.1) {$u=p+\Psi(p)$};
  \node[ccbox, minimum width=2.5cm] (eff) at (3.2,-1.1) {$F_{\mathrm{eff}}(p)=0$};

  \draw[ccarrow] (p) -- node[above, cclabel] {solve shadow} (shadow);
  \draw[ccarrow] (shadow) -- node[above, cclabel] {unique graph} (psi);
  \draw[ccdarrow] (p.south) -- ++(0,-0.55) -| (u.west);
  \draw[ccdarrow] (psi.south) -- ++(0,-0.55) -| (u.east);
  \draw[ccarrow] (u) -- node[above, cclabel] {substitute} (eff);
\end{tikzpicture}%
}
\caption{Exact nonlinear elimination: solve the shadow equation for
$q=\Psi(p)$, then substitute into the observed equation.}
\label{fig:nonlinear-elimination}
\end{figure}

\begin{proposition}[Constructive semilinear shadow solve]
\label{prop:semicolon-shadow-elimination}\lean{semilinear_shadow_elimination}
Let
\[
F(u)=Lu+N(u)-f
\]
on a Hilbert space $\cH$, where $L:\cH\to\cH$ is bounded linear,
$N:\mathcal D\to\cH$ is a nonlinear map, and $f\in\cH$. Fix an orthogonal
projection $P$, set $Q:=I-P$, and let $B_P\subseteq \cH_P$ and
$B_Q\subseteq \cH_Q$ be sets such that $p+q\in\mathcal D$ for every
$p\in B_P$ and $q\in B_Q$, and with $B_Q$ complete in the ambient norm.
Assume:
\begin{enumerate}[label=(\roman*), leftmargin=2em, itemsep=0.2em]
\item $QLQ:\cH_Q\to\cH_Q$ is invertible;
\item for every $p\in B_P$, the map
      \[
      T_p(q):=(QLQ)^{-1}\bigl(Qf-QLP\,p-QN(p+q)\bigr)
      \]
      maps $B_Q$ into itself;
\item there exists $\kappa\in[0,1)$ such that for every $p\in B_P$ and all
      $q_1,q_2\in B_Q$,
      \[
      \|T_p(q_1)-T_p(q_2)\|\le \kappa \|q_1-q_2\|.
      \]
\end{enumerate}
Then:
\begin{enumerate}[label=(\roman*), leftmargin=2em, itemsep=0.2em]
\item for each $p\in B_P$ there exists a unique $\Psi(p)\in B_Q$ satisfying
      \[
      QF\bigl(p+\Psi(p)\bigr)=0;
      \]
\item the map $\Psi:B_P\to B_Q$ is given constructively by the fixed-point
      iteration $\Psi(p)=\lim_{n\to\infty}T_p^{\,n}(q_0)$ for any $q_0\in B_Q$;
\item the full semilinear equation $F(u)=0$ on
      $\{p+q:\,p\in B_P,\ q\in B_Q\}$ is exactly equivalent to the reduced law
      \[
      F_{\mathrm{eff}}(p)
      =
      \CCBlockPP{L}p + \CCBlockPQ{L}\Psi(p) + P N\bigl(p+\Psi(p)\bigr) - Pf
      =0.
      \]
\end{enumerate}
\end{proposition}

(Proof in Appendix~\ref{sec:appendix-elimination}.)

\begin{corollary}[Linear resolvent problem recovered]
\label{cor:linear-shadow-specialization}\lean{linear_shadow_specialization}
Let $F(u)=(I-zA)u-g$ with $A:\cH\to\cH$ linear, $g\in\cH$, and $P$ an
orthogonal projection. Assume $(I-z\CCBlockQQ{A})$ is invertible on
$\cH_{\CCComp}$. Then the unique shadow solve is
\[
\Psi(p)
=
(I-z\CCBlockQQ{A})^{-1}\bigl(Qg+z\CCBlockQP{A}p\bigr),
\]
and the reduced observed equation becomes
\[
\bigl(I-z\CCBlockPP{A}-z^2\CCSigmaP{P}{z}\bigr)p
=
Pg+z\CCBlockPQ{A}(I-z\CCBlockQQ{A})^{-1}Qg.
\]
In particular, if $Qg=0$, then the reduced law is
\[
\bigl(I-z\CCAeff{A}{P}(z)\bigr)p=Pg,
\]
so the linear Feshbach--Schur identity of
Theorem~\ref{thm:feshbach-schur} is recovered exactly.
\end{corollary}

(Proof in Appendix~\ref{sec:appendix-elimination}.)

\subsection{Global elimination and singular regimes}

The preceding proposition gives a constructive local theorem in a contraction
regime. The next results remove that smallness hypothesis on a large global
class and isolate the exact obstruction to further continuation.

\begin{theorem}[Global shadow elimination by strong monotonicity]\lean{global_monotone_shadow}
\label{thm:global-monotone-shadow}
Let $U_P\subseteq \cH_P$ and assume the shadow equation is defined on
$U_P\times \cH_Q$. For each $p\in U_P$, write
\[
G_p(q):=F_Q(p,q).
\]
Assume:
\begin{enumerate}[label=(\roman*), leftmargin=2em, itemsep=0.2em]
\item each $G_p:\cH_Q\to\cH_Q$ is hemicontinuous;
\item there exists $m>0$ such that for every $p\in U_P$ and all
      $q_1,q_2\in\cH_Q$,
      \[
      \operatorname{Re}\langle G_p(q_1)-G_p(q_2),\,q_1-q_2\rangle
      \ge
      m\|q_1-q_2\|^2;
      \]
\item each $G_p$ is coercive:
      \[
      \frac{\operatorname{Re}\langle G_p(q),q\rangle}{\|q\|}
      \longrightarrow +\infty
      \qquad\text{as }\|q\|\to\infty;
      \]
\item there exists $L\ge 0$ such that for every $p_1,p_2\in U_P$ and
      $q\in\cH_Q$,
      \[
      \|G_{p_1}(q)-G_{p_2}(q)\|\le L\|p_1-p_2\|.
      \]
\end{enumerate}
Then:
\begin{enumerate}[label=(\roman*), leftmargin=2em, itemsep=0.2em]
\item for every $p\in U_P$ there exists a unique $\Psi(p)\in\cH_Q$ such that
      $F_Q\bigl(p,\Psi(p)\bigr)=0$;
\item the shadow graph $\Psi:U_P\to\cH_Q$ is globally Lipschitz with
      \[
      \|\Psi(p_1)-\Psi(p_2)\|\le \frac{L}{m}\|p_1-p_2\|;
      \]
\item the full equation $F(u)=0$ on
      $\{\,p+q:\,p\in U_P,\ q\in\cH_Q\,\}$ is exactly equivalent to the
      reduced observed law
      \[
      F_{\mathrm{eff}}(p)=P F\bigl(p+\Psi(p)\bigr)=0.
      \]
\end{enumerate}
No Fredholm assumption is required.
\end{theorem}

(Proof in Appendix~\ref{sec:appendix-elimination}.)

\begin{theorem}[Regular and singular shadow regimes]\lean{regular_singular_shadow}
\label{thm:regular-singular-shadow}
Let $\mathcal O\subseteq \cH_P\times\cH_Q$ be open and suppose
$F_Q:\mathcal O\to\cH_Q$ is $C^1$. Define the regular and singular shadow sets
by
\[
\mathcal R
:=
\{(p,q)\in\mathcal O:\,F_Q(p,q)=0,\ D_qF_Q(p,q)\text{ is invertible}\}
\]
and
\[
\Sigma
:=
\{(p,q)\in\mathcal O:\,F_Q(p,q)=0,\ D_qF_Q(p,q)\text{ is not invertible}\}.
\]
Then every $(p_0,q_0)\in\mathcal R$ admits neighborhoods
$U\subseteq \cH_P$, $V\subseteq \cH_Q$ and a unique $C^1$ map
$\Psi:U\to V$ such that
\[
\Psi(p_0)=q_0
\qquad\text{and}\qquad
F_Q\bigl(p,\Psi(p)\bigr)=0
\quad \text{for all }p\in U.
\]
Moreover,
\[
D\Psi(p)
=
-\bigl[D_qF_Q\bigl(p,\Psi(p)\bigr)\bigr]^{-1}
  D_pF_Q\bigl(p,\Psi(p)\bigr).
\]
Thus $\Sigma$ is exactly the obstruction to local smooth shadow elimination.
\end{theorem}

\begin{proof}
This is the Banach-space implicit function theorem applied to the map
$F_Q$ at a zero $(p_0,q_0)$ with invertible vertical derivative
$D_qF_Q(p_0,q_0)$. The derivative formula follows by differentiating the
identity $F_Q(p,\Psi(p))=0$ and solving for $D\Psi(p)$.
\end{proof}

\begin{proposition}[Gauge-sliced shadow elimination]\lean{gauge_sliced_shadow}
\label{prop:gauge-sliced-shadow}
Assume a group $\Gamma$ acts on $\cH_Q$ by isometries. Let
$S\subseteq \cH_Q$ be a slice with the following properties:
\begin{enumerate}[label=(\roman*), leftmargin=2em, itemsep=0.2em]
\item the shadow equation is gauge-invariant in the sense that
      \[
      F_Q(p,\gamma\cdot q)=0
      \quad\Longleftrightarrow\quad
      F_Q(p,q)=0
      \qquad
      (\gamma\in\Gamma);
      \]
\item every $\Gamma$-orbit contained in the shadow zero set intersects $S$ in
      exactly one point.
\end{enumerate}
Then solving $F_Q(p,q)=0$ modulo gauge is equivalent to solving the restricted
equation
\[
F_Q(p,s)=0
\qquad\text{with } s\in S.
\]
If the restricted equation on $U_P\times S$ satisfies the hypotheses of
Theorem~\ref{thm:global-monotone-shadow} or
Theorem~\ref{thm:regular-singular-shadow}, then the full equation admits a
unique reduced observed law modulo gauge.
\end{proposition}

\begin{proof}
By gauge invariance, the shadow zero set is a union of $\Gamma$-orbits.
Assumption (ii) gives a bijection between those orbits and their unique slice
representatives in $S$. Therefore solving the shadow equation modulo gauge is
the same as solving it on the slice. The final statement is immediate from the
two cited theorems applied to the restricted equation.
\end{proof}

\paragraph{What is now solved, and what remains.}
The open problem of nonlinear elimination is resolved in the theoremic sense
needed by this manuscript. Exact elimination is now available:
\begin{enumerate}[label=(\roman*), leftmargin=2em, itemsep=0.2em]
\item locally by contraction in Proposition~\ref{prop:semicolon-shadow-elimination};
\item globally on a non-Fredholm class by strong monotonicity in
      Theorem~\ref{thm:global-monotone-shadow};
\item near regular singularity-free points by the implicit-function regime of
      Theorem~\ref{thm:regular-singular-shadow};
\item modulo gauge once a slice is fixed, by
      Proposition~\ref{prop:gauge-sliced-shadow}.
\end{enumerate}
What remains open is no longer elimination itself, but model-specific
verification of these hypotheses and branch selection on the singular set
$\Sigma$ when multiple continuations become admissible.
